%------------------------------------------------------------------------------%
% Luis A. D'Afonseca
%
%------------------------------------------------------------------------------%
\documentclass[a4paper,12pt,fleqn]{article}

\usepackage{style_avaliacao}
\usepackage{systeme}

\ShowAnswers

\newcommand{\D}[2]{\frac{\partial {#1}}{\partial {#2}}}
\newcommand{\DS}[2]{\frac{\partial^2 {#1}}{\partial {#2}^2}}
\newcommand{\DM}[3]{\frac{\partial^2 {#1}}{\partial {#2} \partial {#3}}}

%------------------------------------------------------------------------------%
\begin{document}

\makeHeader{IS}{Prova 2 -- Turma 8}{02/08/2024}

% 01
%------------------------------------------------------------------------------%
\exercise{20}
Use substituição simples para calcular a integral
\(
  \int \sqrt{x}\sin\left(x^{\sfrac{3}{2}}+1\right)dx
\)

\begin{answer}
  Queremos calcular
  \[
    F = \int \sqrt{x}\sin\left(x^{\sfrac{3}{2}}+1\right)dx
  \]
  Faremos a substituição
  \[
    u = x^{\sfrac{3}{2}}+1
    \qquad
    du = \frac{3}{2} x^{\sfrac{1}{2}} dx = \frac{3}{2} \sqrt{x} dx
    \qquad
     \sqrt{x} dx = \frac{2}{3} du
  \]
  Assim
  \begin{align*}
    F
    & = \int \sin\left(x^{\sfrac{3}{2}}+1\right) \sqrt{x} dx \\[2mm]
    & = \int \sin\left(u\right) \frac{2}{3} du \\[2mm]
    & = \frac{2}{3} \int \sin\left(u\right) du \\[2mm]
    & = \frac{2}{3} \cos\left(u\right) + C \\[2mm]
    & = \frac{2}{3} \cos\left(x^{\sfrac{3}{2}}+1\right) + C
  \end{align*}
  \clearpage
\end{answer}

% 02
%------------------------------------------------------------------------------%
\exercise{20}
Use integração por partes para calcular a integral
\(
\int_1^e x^3\ln(x) dx
\)

\begin{answer}
  Queremos calcular
  \[
    I = \int_1^e x^3\ln(x) dx
  \]
  Primeiro vamos encontrar a primitiva
  \[
    F = \int x^3\ln(x) dx
  \]
  Usamos
  \[
    u = \ln(x) \qquad du = \frac{1}{x} dx
  \]
  \[
    dv = x^3 dx \qquad v = \frac{x^4}{4}
  \]
  Assim
  \begin{align*}
    F(x)
    & = \int x^3 \ln(x) dx \\[2mm]
    & = \ln(x) \frac{x^4}{4} - \int \frac{x^4}{4} \frac{1}{x} dx \\[2mm]
    & = \frac{x^4\ln(x) }{4} - \frac{1}{4}\int x^3 dx \\[2mm]
    & = \frac{x^4\ln(x) }{4} - \frac{1}{4} \frac{x^4}{4} + C \\[2mm]
    & = \frac{x^4\ln(x) }{4} - \frac{x^4}{16} + C
  \end{align*}
  Agora podemos aplicar o Teorema Fundamental do Cálculo
  \begin{align*}
    I
    & = \int_1^e x^3\ln(x) dx \\[2mm]
    & = F(x)\evalat{1}{e} \\[2mm]
    & = F(e) - F(1) \\[2mm]
    & = \left(\frac{e^4\ln(e) }{4} - \frac{e^4}{16} + C\right)
      - \left(\frac{1^4\ln(1) }{4} - \frac{1^4}{16} + C\right) \\[2mm]
    & = \frac{e^4}{4} - \frac{e^4}{16}
      - \frac{0}{4}   + \frac{1}{16} \\[2mm]
    & = \frac{4e^4 - e^4 + 1}{16} \\[2mm]
    & = \frac{3e^4 + 1}{16}
  \end{align*}
  \clearpage
\end{answer}

% 03
%------------------------------------------------------------------------------%
\exercise{20}
Calcule a integral
\(
  \int \cos^5(x) dx
\)

\begin{answer}
  \begin{align*}
    F
    & = \int \cos^5(x) dx \\[2mm]
    & = \int \cos^4(x) \cos(x) dx \\[2mm]
    & = \int \left(\cos^2(x)\right)^2 \cos(x) dx \\[2mm]
    & = \int \left(1 - \sin^2(x)\right)^2 \cos(x) dx
  \end{align*}
  \[
    u = \sin(x) \qquad du = \cos(x) dx
  \]
  \begin{align*}
    F
    & = \int \left(1 - u^2\right)^2 du \\[2mm]
    & = \int 1 - 2u^2 + u^4 du \\[2mm]
    & = u - 2\frac{u^3}{3} + \frac{u^5}{5} + C \\[2mm]
    & = u - \frac{2}{3}u^3 + \frac{1}{5}u^5 + C\\[2mm]
    & = \sin(x) - \frac{2}{3}\sin^3(x) + \frac{1}{5}\sin^5(x) + C
  \end{align*}
  \clearpage
\end{answer}

% 04
%------------------------------------------------------------------------------%
\exercise{20}
Calcule a integral
\(
  \int \frac{dx}{x^3 + x^2 - 2x}
\)

\begin{answer}
  \[
    F = \int \frac{dx}{x^3 + x^2 - 2x}
  \]
  Vamos aplicar frações parciais, para isso precisamos das raizes de \(Q(x)=x^3 + x^2 - 2x\)
  \begin{align*}
    x^3 + x^2 - 2x & = 0 \\
    x(x^2+x-2) & = 0
  \end{align*}
  Uma das raizes de $Q$ é zero, $x_1=0$, usamos Bhaskara para encontrar as demais
  \[
    x^2+x-2 = 0
  \]
  \[
    \Delta
    = b^2 - 4ac
    = (1)^2 - 4\times (1) \times (-2)
    = 1 + 8
    = 9
  \]
  \[
    \sqrt{\Delta} = 3
  \]
  \[
    x
    = \frac{-b\pm\sqrt{\Delta}}{2a}
    = \frac{1\pm 3}{2}
  \]
  \[
    x_2
    = \frac{1-3}{2}
    = \frac{-2}{2}
    = -1
  \]
  \[
    x_3
    = \frac{1+3}{2}
    = \frac{4}{2}
    = 2
  \]
  Podemos fatorar $Q$
  \[
    Q(x) = x(x+1)(x-2)
  \]
  Usando frações parciais
  \[
    \frac{1}{x^3 + x^2 - 2x} = \frac{A}{x} + \frac{B}{x+1} + \frac{C}{x-2}
  \]
  Multiplicando os dois lados por \(x(x+1)(x-2)\)
  \begin{align*}
      1
      & = A(x+1)(x-2) + Bx(x-2) + Cx(x+1) \\
      & = A(x^2 -2x + x -2) + B(x^2 -2x) + C(x^2+x) \\
      & = Ax^2 -Ax + 2A + Bx^2 -2Bx + Cx^2 + Cx \\
      & = x^2(A+B+C) + x(-A-2B+ C) + 2A
    \end{align*}
  Igualando os polinômios construimos o sistema linear
  \[
  \systeme[ABC]{
       A +B + C = 0,
      -A -2B +C = 0,
      2A        = 1
    }
  \]
  Verificamos que \(A=\sfrac{1}{2}\) e reduzimos o sistema para
  \[
  \systeme[ABC]{
      B +  C =  -\frac{1}{2},
      -2B + C = \frac{1}{2}
    }
  \]
  Isolando $B$ na primeira equação temos $B=\sfrac{1}{2}-C$, substituindo na segunda temos
  \begin{align*}
    -2B + C                          & = \frac{1}{2}    \\[2mm]
    -2\left(\frac{1}{2}-C\right) + C & = \frac{1}{2}    \\[2mm]
    -1 -2C + C                       & = \frac{1}{2}    \\[2mm]
    -C                               & = \frac{1}{2}+ 1 = \frac{3}{2}    \\[2mm]
    C                                & = -\frac{3}{2}
  \end{align*}
  portanto
  \[
    B
    = \frac{1}{2} - C
    = \frac{1}{2} + \frac{3}{2}
    = \frac{4}{2}
    = 2
  \]
  temos então que
  \[
    \frac{1}{{x^3 + x^2 - 2x}}
    = \frac{A}{x} + \frac{B}{x+1} + \frac{C}{x-2}
    = \frac{1}{2x} + \frac{2}{x+1} - \frac{3}{2(x-2)}
  \]
  então
  \begin{align*}
      \int \frac{1}{{x^3 + x^2 - 2x}} dx
      & = \int \frac{1}{2x} + \frac{2}{x+1} - \frac{3}{2(x-2)} dx \\[2mm]
      & = \frac{1}{2} \int \frac{1}{x} dx
        + 2\int \frac{1}{x+1} dx
        - \frac{3}{2} \int \frac{1}{x-2} dx \\[2mm]
      & = \frac{1}{2}\ln\abs{x} + 2\ln\abs{x+1} - \frac{3}{2}\ln\abs{x-2} + C
    \end{align*}
  \clearpage
\end{answer}


%\begin{answer}
%  \[
%  F = \int \frac{2x^2 -3 x + 3}{x^3-2x^2-3x} dx
%  \]
%  Vamos aplicar frações parciais, para isso precisamos das raizes de \(Q(x)=x^3-2x^2-3x\)
%  \begin{align*}
%    x^3-2x^2-3x & = 0 \\
%    x(x^2-2x-3) & = 0
%  \end{align*}
%  Uma das raizes de $Q$ é zero, $x_1=0$, usamos Bhaskara para encontrar as demais
%  \[
%  x^2-2x-3 = 0
%  \]
%  \[
%  \Delta
%  = b^2 - 4ac
%  = (-2)^2 - 4\times 1 \times (-3)
%  = 4 + 12
%  = 16
%  \]
%  \[
%  \sqrt{\Delta} = 4
%  \]
%  \[
%  x
%  = \frac{-b\pm\sqrt{\Delta}}{2a}
%  = \frac{2\pm 4}{2}
%  \]
%  \[
%  x_2
%  = \frac{2-4}{2}
%  = \frac{-2}{2}
%  = -1
%  \]
%  \[
%  x_3
%  = \frac{2+4}{2}
%  = \frac{6}{2}
%  = 3
%  \]
%  Podemos fatorar $Q$
%  \[
%  Q(x) = x(x-3)(x+1)
%  \]
%  Usando frações parciais
%  \[
%  \frac{2x^2 -3 x + 3}{x^3-2x^2-3x} = \frac{A}{x} + \frac{B}{x-3} + \frac{C}{x+1}
%  \]
%  Multiplicando os dois lados por \(x(x-3)(x+1)\)
%  \begin{align*}
%    2x^2 -3x + 3
%    & = A(x-3)(x+1) + Bx(x+1) + C x(x-3) \\
%    & = A(x^2 +x -3x -3) + B(x^2+x) + C(x^2-3x) \\
%    & = Ax^2 - 2Ax -3A + Bx^2 +Bx +Cx^2 - 3Cx \\
%    & = (A+B+C)x^2 + (-2A+B-3C)x -3A
%  \end{align*}
%  Igualando os polinômios construimos o sistema linear
%  \[
%  \systeme[ABC]{
%    A +B + C = 2,
%    -2A +B -3C =-3,
%    -3A        = 3
%  }
%  \]
%  Verificamos que \(A=-1\) e reduzimos o sistema para
%  \[
%  \systeme[ABC]{
%    B +  C =  3,
%    B - 3C = -5
%  }
%  \]
%  Isolando $B$ na primeira equação temos $B=3-C$, substituindo na segunda temos
%  \begin{align*}
%    B - 3C     & = -5 \\
%    (3-C) - 3C & = -5 \\
%    3 - C - 3C & = -5 \\
%    -4C        & = -8 \\
%    C          & = 2
%  \end{align*}
%  portanto $B=1$, temos então que
%  \[
%  \frac{2x^2 -3 x + 3}{x^3-2x^2-3x} = -\frac{1}{x} + \frac{1}{x-3} + \frac{2}{x+1}
%  \]
%  então
%  \begin{align*}
%    \int \frac{2x^2 -3 x + 3}{x^3-2x^2-3x} dx
%    & = -\int \frac{1}{x} dx + \int \frac{1}{x-3} dx + 2\int \frac{1}{x+1} dx \\[2mm]
%    & = -\ln\abs{x} + \ln\abs{x-3} + 2\ln\abs{x+1} + C \\[2mm]
%    & = \ln\abs{x-3} - \ln\abs{x} + 2\ln\abs{x+1} + C
%  \end{align*}
%  \clearpage
%\end{answer}

% 05
%------------------------------------------------------------------------------%
\exercise{20}
Calcule a limite da sequência \(a_k = \frac{3^{k+1}}{5^k}\)

\begin{answer}
  \begin{align*}
    \lim a_k
    & = \lim \frac{3^{k+1}}{5^k} \\[2mm]
    & = \lim 3 \frac{3^{k}}{5^k} \\[2mm]
    & = 3\lim \left(\frac{3}{5}\right)^k \\[2mm]
    & = 3 \times 0 \\[2mm]
    & = 0
  \end{align*}
  \clearpage
\end{answer}

%------------------------------------------------------------------------------%
\end{document}
%------------------------------------------------------------------------------%
